\section{Remaining Work}

There is still significant potential to expand the breadth of our project over the second half of the semester. So far we have tested mostly uniform perturbations we plan to augment this to test many different types of affine transforms that the method may encounter such as rotations and translations. We also want to look at domain over-sizing, the domain bounding box discussed previously defines the region in which to begin subdividing into octants. If this box is significantly larger than the true domain of the particles it will cause the octree to have large regions of empty space which requires the octree to become deeper thus inducing inefficiencies. A study on the sizing of the bounding box against the build time should be conducted. An interesting parameter to study would be $L$ which is the maximal depth an octree can be resolved to. The two values usually used are 32-bit integers which allow resolution to a depth of 10 levels and 64-bit integers which allow resolution to a depth of 21. The tradeoff can sacrifice some accuracy but can reduce the amount of memory movement significantly. This is especially of interest since the radix sort and gather kernels are bottlenecks in the local case both of which are memory bound kernels which can be improved by reducing global memory pressure.

We have mostly looked at NSight Systems results to understand communication bottlenecks and which sections of code are the largest bottlenecks, assuming it is of interest we will likely look at some NSight Compute results for the kernels to determine specifically which instructions are causing issues. The radix sort and gather kernels are both from the NVIDIA Core Compute Library (CCCL) and should be very optimized, so optimizing the kernel code is unlikely to provide significant speedup. This means that we should look into stream and graph based execution of the these operations although another optimization idea would be kernel fusion to reduce data movement and the overhead of multiple kernel launches. These two studies would likely be interesting and both options may be explored in the final report.
